\section{Instrumentation Depth and Performance}

We measured the instrumentation depth at startup using \texttt{AFL\_DEBUG=1}. To measure the complete library, we linked a minimal driver against \texttt{libpng14.a} using \texttt{-Wl,--whole-archive}, yielding 14,316 instrumented edges. In contrast, our final harness binary (\texttt{bin/png\_fuzz}) reported 10,271 edges. The gap comes from how the two binaries were linked. With
\texttt{--whole-archive}, the measurement binary includes the whole instrumented libpng archive. The final \texttt{bin/png\_fuzz} binary, however, was linked normally, so it only includes the libpng parts actually needed by the harness. This leaves out unused code, such as encoder-related functions. The harness adds a few extra instrumented edges, but not enough to compensate for the libpng code left out.

For the ASan fork-mode campaign using \texttt{bin/png\_fuzz}, AFL++ reported 3158 reached edges and a map density of 30.73\%. This directly correlates with our harness count (3,158 out of 10,271 equals roughly 30.75\%, calculated against the 10,276 target map size). This means that the campaign reached roughly one third of the instrumented edge surface in the final fuzzing binary: the remaining edges represent unexplored code path. Because PNG is highly structured, reaching deeper branches requires precise combinations of chunk types, bit depths, and error states that the fuzzer simply did not trigger during the campaign.

Finally, we compared three execution speeds. The no-sanitizer fork-mode baseline hit 692.06 exec/s, while adding ASan dropped it to 310.50 exec/s (a 2.23x slowdown) due to computationally heavy memory checks and shadow-memory tracking. However, ASan in persistent mode reached 6,199.60 exec/s (a 20x speedup).  While fork mode forces the OS to create a new process and execute file I/O for every input, persistent mode completely bypasses this bottleneck by continuously feeding inputs to a single, long-lived target process via a tight \texttt{\_\_AFL\_LOOP}.